\documentclass[12pt]{article}
%---DOCUMENT MARGINS---
\usepackage{geometry} % Required for adjusting page dimensions and margins
\geometry{
	paper=a4paper, % Paper size, change to letterpaper for US letter size
	top=4cm, % Top margin
	bottom=2cm, % Bottom margin
	left=2cm, % Left margin
	right=2cm, % Right margin
	headheight=3cm, % Header height
	%footskip=1.5cm, % Space from the bottom margin to the baseline of the footer
	headsep=0.5cm, % Space from the top margin to the baseline of the header
	%showframe, % Uncomment to show how the type block is set on the page
}
\usepackage{cmap}

\usepackage[T1,T2A]{fontenc}
\usepackage{fontspec}
\setmainfont[
	BoldFont={* Bold},
	ItalicFont={* Italic},
	BoldItalicFont={* BoldItalic}
]{DejaVu Serif Condensed}
\setsansfont{DejaVu Sans}
\setmonofont{Dejavu Sans Mono}
\usepackage[math-style=TeX]{unicode-math}
\setmathfont{Latin Modern Math}

\usepackage[nobottomtitles*]{titlesec}
\titleformat
{\section} % command
{\fontsize{14}{14}\sffamily\bfseries} % format
{} % label
{0pt} % sep
{} % before-code
\titlespacing{\section}{0pt}{0.75em}{0.25em}
\newcommand{\tl}{}
\newcommand{\ml}{}
\newcommand{\problem}[1]{%
	\section[#1]{#1 \fontsize{14}{14}\selectfont{\hfill{\normalfont{
				\emoji{hourglass-not-done}\tl \space\space \emoji{floppy-disk} \ml}}}}
}
\AddToHook{cmd/section/before}{\clearpage}

\usepackage[dvipsnames]{xcolor}
\titleformat
{\subsection} % command
{\fontsize{14}{14}\sffamily\bfseries} % format
{\includesvg[width=0.035\textwidth, inkscapelatex=false]{lion}\space} % label
{0pt} % sep
{} % before-code
\titlespacing{\subsection}{0pt}{0.75em}{0.25em}

\setlength{\parskip}{0.5em}
\setlength{\parindent}{0pt}
\renewcommand{\baselinestretch}{1.2}
\sloppy

\usepackage{listings}
\definecolor{lstbg}{RGB}{250,250,252}
\definecolor{lstframe}{RGB}{220,220,225}
\lstdefinestyle{mystyle}{
	backgroundcolor=\color{lstbg},
	frame=single,
	rulecolor=\color{lstframe},
	basicstyle=\ttfamily\small,
	keywordstyle=\bfseries,
	breaklines=true,
	aboveskip=0.5\baselineskip,
	belowskip=0\baselineskip  
}
\lstset{style=mystyle, language=C++}

\usepackage{fancyhdr}
\pagestyle{fancy}
\setlength{\headwidth}{\textwidth}
\newcommand{\round}{}
\newcommand{\problemName}{}
\newcommand{\lang}{English}
\fancyhead[L]{
	\begin{minipage}{0.4\textwidth}
		\bf{
			EJOI 2025 \round\\
			\problemName\\
			\emoji{globe-with-meridians} \lang
		}
	\end{minipage}
	\vspace{0.5cm}
}
\fancyhead[R]{%
	\begin{minipage}{0.6\textwidth}
		\includesvg[width=\textwidth, inkscapelatex=false]{logo}
	\end{minipage}
	\vspace{0.5cm}
}
\usepackage{lastpage}
\fancyfoot[C]{\problemName\space(\lang) \thepage\ / \pageref*{LastPage}}

\raggedbottom

\usepackage{amsmath}
\usepackage{stmaryrd}

\usepackage{graphicx}
\graphicspath{{./}}
\usepackage[export]{adjustbox}
\usepackage{wrapfig}
\makeatletter
\patchcmd\WF@putfigmaybe{\lower\intextsep}{}{}{\fail}
\AddToHook{env/wrapfigure/begin}{\setlength{\intextsep}{0pt}}
\makeatother
\usepackage[inkscapearea=page, inkscapepath=./svg-inkscape]{svg}
\svgpath{{./}}

\usepackage{makecell}
\usepackage{tabularray}
\AtBeginEnvironment{table}{\vspace{-0.2cm}}
\AtEndEnvironment{table}{\vspace{-0.2cm}}
\usepackage{float}
\usepackage{placeins}
\usepackage{caption}
\captionsetup[table]{
	skip=0.25em,
	singlelinecheck=false, justification=justified
}

\usepackage{enumitem}
\setlist{itemsep=-0.4em, leftmargin=22pt, topsep=-\parskip}
\AfterEndEnvironment{itemize}{\vspace{\parskip}}
\newcommand{\tabitem}{\indent~~\llap{\textbullet}~~}

\usepackage{hyperref}
\hypersetup{
	colorlinks=true,
	citecolor=blue,
	linkcolor=blue,
	urlcolor=cyan,
}

\usepackage{emoji}

\usepackage[english]{babel}

\begin{document}

\renewcommand{\round}{Dan 1}
\renewcommand{\problemName}{Zadatak Zatvor}
\renewcommand{\lang}{Hrvatski}

\renewcommand{\tl}{$2$ sek.}
\renewcommand{\ml}{$1024$ MB}
\problem{\problemName}

Alice i Bob su nepravedno osuđeni na zatvor s maksimalnim osiguranjem. Sada moraju isplanirati bijeg. Da bi to učinili, moraju moći komunicirati što učinkovitije (posebno Alice mora slati dnevne informacije Bobu). Međutim, ne mogu se sastati i mogu razmjenjivati informacije samo putem bilješki napisanih na salvetama. Svaki dan Alice želi poslati novu informaciju Bobu - broj između $0$ i $N-1$. Za svaki ručak Alice dobiva tri salvete i na svaku salvetu napiše broj između $0$ i $M-1$ (mogu se ponavljati) i ostavlja ih na svom sjedalu. Zatim njihov neprijatelj, Charly, uništava jednu od salveta i pomiješa preostale dvije. Konačno, Bob pronalazi dvije preostale salvete i čita brojeve na njima. Mora točno dekodirati izvorni broj koji mu je Alice htjela poslati. Na salvetama je ograničen prostor, pa je $M$ fiksan. Međutim, Alicein i Bobov cilj je maksimizirati protok informacija, tako da mogu slobodno odabrati što veći broj $N$. Pomozi Alici i Bobu tako što ćeš za svakog od njih primijeniti strategiju, pokušavajući maksimizirati vrijednost od $N$.

\subsection{Implementacijski detalji}
Budući da se radi o komunikacijskom problemu, vaš će se program izvršavati u dva odvojena izvršavanja (jedno za Alice i jedno za Boba) koja ne mogu dijeliti podatke ili komunicirati na bilo koji drugi način osim onog opisanog ovdje. Potrebno je implementirati tri funkcije:

\begin{lstlisting}
 int setup(int M);
\end{lstlisting}

Ova funkcija će biti pozvana jednom na početku Alicinog izvršavanja vašeg programa i jednom na početku Bobovog izvršavanja. Dodijeljen je $M$ i funkcija mora vratiti željeni $N$. Oba poziva funkcije \texttt{setup} moraju vratiti isti $N$.

\begin{lstlisting}
 std::vector<int> encode(int A);
\end{lstlisting}

Ova funkcija implementira Alisinu strategiju. Bit će pozvana s brojem za kodiranje $A$ ($0 \leq A < N$) i mora vratiti tri broja $W_1, W_2, W_3$ ($0 \leq W_i < M$) koji kodiraju $A$. Ova funkcija bit će pozvana ukupno $T$ puta -- jednom dnevno (vrijednosti $A$ mogu se ponavljati između dana).

\begin{lstlisting}
 int decode(int X, int Y);
\end{lstlisting}

Ova funkcia implementira Bobovu strategiju. Bit će pozvana s dva od tri broja koje vraća \texttt{encode} nekim redoslijedom. Mora vratiti istu vrijednost $A$ koju je \texttt{encode} primio. Ova će se funkcija također pozivati $T$ puta -- što odgovara $T$ pozivima funkcije \texttt{encode}; bit će istim redoslijedom. Svi pozivi funkcije \texttt{encode} dogodit će se prije svih poziva funkcije \texttt{decode}.

\subsection{Ograničenja}

\begin{itemize}
\item $M \leq 4300$
\item $T = 5000$
\end{itemize}

\subsection{Bodovanje}

Za određeni podzadatak, udio S bodova koje dobijete ovisi o najmanjem $N$ koji vraća \texttt{setup} na bilo kojem testu u tom podzadatku. Također ovisi o $N^*$, što je ciljana vrijednost od $N$ koja vam je potrebna da biste dobili puni broj bodova za podzadatak:

\begin{itemize}
\item Ako je vaše rješenje krivo na bilo kojem primjeru, tada je $S = 0$.
\item Ako je $N \geq N^*$, tada je $S = 1.0$.
\item Ako je $N < N^*$, tada je $S = \max\left(0.35 \max\left(\frac{\log(N) - 0.985 \log(M)}{\log(N^*) - 0.985 \log(M)}, 0.0\right)^{0.3} + 0.65 \left(\frac{N}{N^*}\right)^{2.4}, 0.01\right)$.
\end{itemize}

\subsection{Podzadaci}

\begin{table}[H]
\begin{tblr}{|Q[c,m]|Q[c,m]|Q[c,m]|Q[c,m]|}
\hline
\textbf{Podzadatak} & \textbf{Bodovi} & $M$ & $N^*$ \\
\hline
$1$ & $10$ & $700$ & $82017$ \\
\hline
$2$ & $10$ & $1100$ & $202217$ \\
\hline
$3$ & $10$ & $1500$ & $375751$ \\
\hline
$4$ & $10$ & $1900$ & $602617$ \\
\hline
$5$ & $10$ & $2300$ & $882817$ \\
\hline
$6$ & $10$ & $2700$ & $1216351$ \\
\hline
$7$ & $10$ & $3100$ & $1603217$ \\
\hline
$8$ & $10$ & $3500$ & $2043417$ \\
\hline
$9$ & $10$ & $3900$ & $2536951$ \\
\hline
$10$ & $10$ & $4300$ & $3083817$ \\
\hline
\end{tblr}
\end{table}
\FloatBarrier

\subsection{Primjer}

Razmotrimo sljedeći primjer s $T = 5$. Ovdje imamo shemu kodiranja gdje Alice šalje tri jednaka broja za kodiranje $0$ ili tri različita broja za kodiranje $1$. Primijetite da Bob može dekodirati izvorni broj iz bilo koja dva od tri broja koje je Alice poslala.

\begin{table}[H]
\begin{tblr}{|Q[c,m]|Q[c,m]|Q[c,m]|}
\hline
\textbf{\makecell{Osoba}} & \textbf{\makecell{Poziv funkcije}} & \textbf{Povratna vrijednost} \\
\hline
Alice & \texttt{setup(10)} & \texttt{2} \\
\hline
Bob & \texttt{setup(10)} & \texttt{2} \\
\hline
Alice & \texttt{encode(0)} & \texttt{\{5, 5, 5\}} \\
\hline
Alice & \texttt{encode(1)} & \texttt{\{8, 3, 7\}} \\
\hline
Alice & \texttt{encode(1)} & \texttt{\{0, 3, 1\}} \\
\hline
Alice & \texttt{encode(0)} & \texttt{\{7, 7, 7\}} \\
\hline
Alice & \texttt{encode(1)} & \texttt{\{6, 2, 0\}} \\
\hline
Bob & \texttt{decode(5, 5)} & \texttt{0} \\
\hline
Bob & \texttt{decode(8, 7)} & \texttt{1} \\
\hline
Bob & \texttt{decode(3, 0)} & \texttt{1} \\
\hline
Bob & \texttt{decode(7, 7)} & \texttt{0} \\
\hline
Bob & \texttt{decode(2, 0)} & \texttt{1} \\
\hline
\end{tblr}
\end{table}

\subsection{Ocjenjivač probnih primjera}

Za primjer ocjenjivača, svi pozivi na \texttt{encode} i \texttt{decode} bit će u istom izvršavanju vašeg programa. Osim toga, \texttt{setup} će se pozvati samo jednom (za razliku od dva puta, jednom po izvršavanju, kao u sustavu ocjenjivanja).

Ulaz je samo jedan cijeli broj -- $M$. Zatim će ispisati $N$ koje je vaš \texttt{setup} vratio. Zatim će pozvati funkcije \texttt{encode} i \texttt{decode} ovim redoslijedom $T$ puta sa slučajno generiranim brojevima od $0$ do $N-1$ i slučajno generiranim izborima koja dva od tri broja iz \texttt{encode} dati \texttt{decode} (i kojim redoslijedom). Ispisat će poruku o pogrešci ako vaše rješenje nije uspjelo.

\end{document}
